%----------------------------------------------------------------------------------------
% TOEPASSINGEN VAN AFGELEIDEN. Alle opgaven zijn eigen werk.
%
% Opbouw: Start, leerstof (theorie, werkwijze of uitleg, voorbeeld, korte oefening), en op een
% nieuwe pagina de Opdrachten (door elkaar, stijgende moeilijkheid), Koppel en Test jezelf.
% Labels van dit hoofdstuk beginnen met h4-.
%----------------------------------------------------------------------------------------

\mychapter{4}{Toepassingen van afgeleiden}

\lkr{Leerplan B+S': LPD 19-22 (afgeleiden), het deel over het verloop van functies (stijgen, dalen, extrema, holheid, buigpunten) en optimaliseringsproblemen. Over \textit{hol} en \textit{bol} bestaan verschillende afspraken; dit boek koppelt ze telkens aan de vorm van een dal- of bergparabool, zodat er geen misverstand kan zijn. De stelling dat $f' > 0$ stijgen geeft, wordt hier niet bewezen (ze volgt uit de middelwaardestelling).}

In Algemene functies las je toppen en dalen af op een grafiek. Nu kun je ze berekenen: in een top of een dal is de raaklijn horizontaal, dus de afgeleide nul. Daarmee los je ook praktische vragen op: hoe maak je een blik met zo weinig mogelijk metaal, of een weide zo groot mogelijk?

\begin{oefeningen}[Start]
\oef{h4-start-1}

\oef{h4-verken}
\end{oefeningen}

%========================================================================================
\section{Stijgen, dalen en extrema}

\begin{theorie}
Voor een functie $f$ die afleidbaar is op een interval $I$:
\begin{itemize}
\item is $\formule{f'(x) > 0}$ op $I$, dan is $f$ \textbf{\textcolor{accent}{stijgend}} op $I$; is $\formule{f'(x) < 0}$, dan is $f$ \textbf{\textcolor{accent}{dalend}};
\item heeft $f$ een extremum in een punt $a$ binnen $I$, dan is $\formule{f'(a) = 0}$;
\item wisselt $f'$ in $a$ van $+$ naar $-$, dan is $f(a)$ een \textbf{\textcolor{accent}{maximum}}; van $-$ naar $+$ een \textbf{\textcolor{accent}{minimum}}.
\end{itemize}
\end{theorie}

Een nulwaarde van $f'$ is niet altijd een extremum: bij $x^3$ is $f'(0) = 0$, maar $f'$ blijft positief en de functie stijgt door. Kijk dus altijd naar het teken van $f'$ links en rechts.

\voorbeeld{}{Bepaal het verloop en de extrema van $f(x) = x^3 - 3x$.\label{vb:h4-extrema}}{%
\naast{%
\uitleg{Leid af, ontbind, en stel de tekentabel van $f'$ op. Zet eronder pijlen voor $f$.}
\begin{stappenvrij}{4.2cm}{3.3cm}
f'(x) = 3x^2 - 3 & \\
= 3(x - 1)(x + 1) & Ontbinden. \\
\text{max. } f(-1) = 2 & $f'$: van $+$ naar $-$. \\
\text{min. } f(1) = -2 & $f'$: van $-$ naar $+$.
\end{stappenvrij}
\figuren{\tekentabel{cccccc}{x & & -1 & & 1 & \\\hline f'(x) & + & 0 & - & 0 & +}{f(x) & \nearrow & 2 & \searrow & -2 & \nearrow}}}{%
\begin{tikzpicture}[x=.6cm,y=.42cm,font=\scriptsize]
  \rooster{-2}{1}{2}{-3}{1}{3}
  \kromme{-2.1}{2.1}{(\x)^3-3*\x}
  \vol{(-1,2)} \vol{(1,-2)}
\end{tikzpicture}}}

\begin{werkwijze}
Bepaal het domein en leid af. Zoek de nulwaarden van $f'$ (en de punten waar $f'$ niet bestaat). Stel een tekentabel op voor $f'$, met daaronder pijlen voor $f$. Een wissel van $+$ naar $-$ is een maximum, van $-$ naar $+$ een minimum. Bereken de functiewaarden.
\end{werkwijze}

\begin{oefeningen}
\oef{h4-tekentabel}
\end{oefeningen}

\voorbeeld{}{Bepaal de extrema van $f(x) = x\,\e^{-x}$.\label{vb:h4-extrema2}}{%
\uitleg{De afgeleide uit hoofdstuk 3. De factor $\e^{-x}$ is altijd positief: het teken hangt alleen af van $1 - x$.}
\begin{stappenbreed}
f'(x) = (1 - x)\e^{-x} & Productregel en kettingregel. \\
f'(x) > 0 \text{ voor } x < 1,\ f'(x) < 0 \text{ voor } x > 1 & \\
\text{max. } f(1) = \tfrac{1}{\e} \approx 0{,}37 & Het enige extremum.
\end{stappenbreed}}

\begin{oefeningen}
\oef{h4-extrema}

\oef{h4-extremaexp}
\end{oefeningen}

%========================================================================================
\section{Holheid en buigpunten}

De afgeleide van $f'$ heet de \textbf{tweede afgeleide} $f''$. Ze zegt hoe de helling verandert: neemt de helling toe, dan buigt de grafiek naar boven.

\begin{theorie}
Voor een functie $f$ die twee keer afleidbaar is op een interval $I$:
\begin{itemize}
\item is $\formule{f''(x) > 0}$ op $I$, dan is de grafiek \textbf{\textcolor{accent}{hol}} op $I$: ze buigt naar boven, zoals een dalparabool ($\cup$);
\item is $\formule{f''(x) < 0}$, dan is de grafiek \textbf{\textcolor{accent}{bol}}: ze buigt naar onder, zoals een bergparabool ($\cap$);
\item een punt waar $f''$ van teken wisselt, is een \textbf{\textcolor{accent}{buigpunt}}.
\end{itemize}
\end{theorie}

\voorbeeld{}{Bepaal de extrema, de holheid en de buigpunten van $f(x) = x^4 - 4x^3$.\label{vb:h4-buigpunt}}{%
\naast{%
\uitleg{Stel twee tekentabellen op: voor $f'$ en voor $f''$. In $0$ is $f'$ nul zonder van teken te wisselen: daar is geen extremum.}
\begin{stappenvrij}{4.3cm}{3.2cm}
f'(x) = 4x^3 - 12x^2 & \\
= 4x^2(x - 3) & \\
\text{min. } f(3) = -27 & Geen extremum in $0$. \\
f''(x) = 12x^2 - 24x & \\
= 12x(x - 2) & \\
\text{buigpunten } (0, 0),\ (2, -16) & $f''$ wisselt van teken.
\end{stappenvrij}
\figuren{\tekentabel{cccccccc}{x & & 0 & & 2 & & 3 & \\\hline f'(x) & - & 0 & - & - & - & 0 & + \\ f''(x) & + & 0 & - & 0 & + & + & +}{f(x) & \cup & 0 & \cap & -16 & \cup & -27 & \cup}}}{%
\begin{tikzpicture}[x=.55cm,y=.075cm,font=\scriptsize]
  \draw[black!14,very thin,xstep=1,ystep=10] (-1,-30) grid (4,20);
  \draw[->,grijs] (-1.3,0) -- (4.4,0) node[right,text=black] {$x$};
  \draw[->,grijs] (0,-31) -- (0,22) node[above,text=black] {$y$};
  \foreach \t in {-1,1,2,3,4}{\node[below,inner sep=1.5pt] at (\t,0) {$\t$};}
  \foreach \t in {-20,-10,10}{\node[left,inner sep=1.5pt] at (0,\t) {$\t$};}
  \kromme{-1.2}{4.25}{(\x)^4-4*(\x)^3}
  \vol{(3,-27)} \vol{(0,0)} \vol{(2,-16)}
\end{tikzpicture}}}

\begin{oefeningen}
\oef{h4-holbol}

\oef{h4-buigpunt}

\oef{h4-buigpunt2}
\end{oefeningen}

%========================================================================================
\section{Optimaliseren}

\voorbeeld{}{Een boer zet een rechthoekige weide af langs een rivier. Aan de rivierkant komt geen hek; voor de drie andere zijden heeft hij $100$ m draad. Hoe groot kan de weide hoogstens worden?\label{vb:h4-weide}}{%
\naast{%
\uitleg{Kies de breedte $x$ loodrecht op de rivier. Dan is de lengte langs de rivier $100 - 2x$.}
\begin{stappenvrij}{4.6cm}{3.2cm}
A(x) = x(100 - 2x) & Oppervlakte. \\
0 < x < 50 & Domein. \\
A'(x) = 100 - 4x & \\
A'(x) = 0 \iff x = 25 & \\
A' > 0 \text{ voor } x < 25 & \\
A' < 0 \text{ voor } x > 25 & Een maximum. \\
A(25) = 25 \cdot 50 = 1250 & In m$^2$.
\end{stappenvrij}}{%
\begin{tikzpicture}[scale=.5,font=\scriptsize]
  \fill[accent!10] (0,0) rectangle (5,2.5);
  \draw[accent,thick] (0,0) -- (0,2.5) -- (5,2.5) -- (5,0);
  \draw[grijs] plot[domain=-.8:5.8,samples=80,smooth] (\x,{.12*sin(360*\x/1.1)});
  \node[below] at (2.5,0) {rivier};
  \node[left] at (0,1.25) {$x$};
  \node[above] at (2.5,2.5) {$100 - 2x$};
\end{tikzpicture}}}

\begin{werkwijze}
Kies een veranderlijke en druk de grootheid die optimaal moet zijn uit als functie van die ene veranderlijke (gebruik de voorwaarde om andere grootheden te elimineren). Bepaal het domein in de context. Leid af, zoek de nulwaarden, en controleer met een tekentabel (of de randen van het domein) dat het een maximum of minimum is. Beantwoord de vraag in de context.
\end{werkwijze}

\begin{oefeningen}
\oef{h4-weide2}
\end{oefeningen}

\voorbeeld{}{Een cilindervormig blik moet $1$ liter ($1000$ cm$^3$) bevatten. Welke afmetingen vragen het minste metaal?\label{vb:h4-blik}}{%
\uitleg{Minimaliseer de totale oppervlakte (bodem, deksel en mantel). Druk de hoogte $h$ uit met de inhoud: $\pi r^2 h = 1000$.}
\begin{stappenbreed}
h = \frac{1000}{\pi r^2},\quad A(r) = 2\pi r^2 + 2\pi r h = 2\pi r^2 + \frac{2000}{r} & $r > 0$. \\
A'(r) = 4\pi r - \frac{2000}{r^2} = 0 \iff r^3 = \frac{500}{\pi} & \\
r = \sqrt[3]{\frac{500}{\pi}} \approx 5{,}42,\quad h = \frac{1000}{\pi r^2} \approx 10{,}84 & $A'$ wisselt van $-$ naar $+$: minimum.
\end{stappenbreed}
Het zuinigste blik is even hoog als breed: $h = 2r$. Want uit $r^3 = \frac{500}{\pi}$ volgt $h = \frac{1000}{\pi r^2} = \frac{2r^3}{r^2} = 2r$.}

\begin{oefeningen}
\oef{h4-optimaal}
\end{oefeningen}

%========================================================================================
\opdrachtenkop

\begin{oefeningen}
\needspace{16\baselineskip}
\oef{h4-o10}

\oef{h4-o2}

\needspace{10\baselineskip}
\oef{h4-o16}

\oef{h4-o12}

\oef{h4-o1}

\needspace{14\baselineskip}
\oef{h4-o15}

\needspace{14\baselineskip}
\oef{h4-o11}

\oef{h4-o3}

\oef{h4-o13}

\oef{h4-o4}

\oef{h4-o14}

\needspace{12\baselineskip}
\oef{h4-o23}

\oef{h4-o19}

\oef{h4-o17}

\oef{h4-o5}

\oef{h4-o20}

\needspace{14\baselineskip}
\oef{h4-o24}

\oef{h4-blik2}

\oef{h4-o21}

\oef{h4-o6}

\oef{h4-o18}

\needspace{10\baselineskip}
\oef{h4-o22}

\oef{h4-o7}

\oef{h4-o26}

\oef{h4-o25}

\oef{h4-o8}

\oef{h4-o27}

\oef{h4-o9}
\end{oefeningen}

%----------------------------------------------------------------------------------------
% KOPPEL: verbanden met andere leerstof
%----------------------------------------------------------------------------------------
\needspace{20\baselineskip}
\opdrachtsectie{Koppel}

\lkr{Statistiek, Rationale functies en Exponentiële en logaritmische functies: kijk na wat al gezien is. Het pakket legt de volgorde niet vast.}

\begin{oefeningen}
\oef{h4-koppel-1}

\oef{h4-koppel-2}

\oef{h4-koppel-3}
\end{oefeningen}

%----------------------------------------------------------------------------------------
\opdrachtsectie{Test jezelf}

\begin{oefeningen}
\oef{h4-t1}

\oef{h4-t2}

\oef{h4-t3}

\oef{h4-t4}

\oef{h4-t5}

\oef{h4-t6}

\oef{h4-t7}
\end{oefeningen}

\herhaalschema{%
  \fv{oef:h4-t1}{vb:h4-extrema}{oef:h4-tekentabel}%
  \fv{oef:h4-t2}{vb:h4-extrema2}{oef:h4-extremaexp}%
  \fv{oef:h4-t3}{vb:h4-buigpunt}{oef:h4-buigpunt}%
  \fv{oef:h4-t4}{vb:h4-weide}{oef:h4-weide2}%
  \fzv{oef:h4-t5}{oef:h4-o2}%
  \fv{oef:h4-t6}{vb:h4-buigpunt}{oef:h4-o10}%
  \fzv{oef:h4-t7}{oef:h4-o22}}

\ifzelfstudie
\section*{Maak je eigen samenvatting}
Sluit dit boek en schrijf in je schrift: het verband tussen het teken van $f'$ en stijgen of dalen, wanneer een nulwaarde van $f'$ een maximum of minimum is, wat $f''$ zegt over hol en bol en wat een buigpunt is, en de werkwijze voor optimaliseringsproblemen. Vergelijk daarna met de kaders.
\fi
